%! Author = chtompki
%! Date = 1/14/26
\documentclass[10pt]{amsart}
% Default margins are too wide all the way around. I reset them here
\setlength{\hoffset}{-.75in}
\setlength{\textwidth}{6.5in}
\setlength{\voffset}{-.5in}
\setlength{\textheight}{9.0in}
\usepackage{amsmath}
\usepackage{amssymb}
\usepackage{amsthm}
\usepackage{lmodern}
\usepackage{microtype}
\usepackage{enumitem}
\usepackage{setspace}
\usepackage{graphicx}
\usepackage{hyperref}
\usepackage{tikz}
\usepackage{listings}
\usepackage{pgfplots}
\pgfplotsset{compat=1.18}
\usepgfplotslibrary{groupplots}


\newenvironment{solution}
{\begin{proof}[Solution]}
{\end{proof}}

\NewDocumentCommand{\vitem}{v}{%
    \item[\texttt{#1}]%
}

\newcommand{\slopeelement}[3]{%
    \pgfmathsetmacro{\slopeDX}{0.14/sqrt(1+(#3)^2)}%
    \pgfmathsetmacro{\slopeDY}{(#3)*\slopeDX}%
    \pgfmathsetmacro{\slopeXA}{(#1)-\slopeDX}%
    \pgfmathsetmacro{\slopeYA}{(#2)-\slopeDY}%
    \pgfmathsetmacro{\slopeXB}{(#1)+\slopeDX}%
    \pgfmathsetmacro{\slopeYB}{(#2)+\slopeDY}%
    \edef\slopecommand{%
        \noexpand\draw[gray,thin]
        (axis cs:\slopeXA,\slopeYA)
        --
        (axis cs:\slopeXB,\slopeYB);%
    }%
    \slopecommand
}

\title{Dynamical Systems\\Midterm Preparation}
\author{Rob Tompkins}
\date{\today}

\begin{document}
    \maketitle
    \date{\today}

\section*{Exam Practice Questions (Ch 1--6.4)}
% Transcribed from M535MTproblems.pdf (MATH 535, Fall 2024).
\noindent\textbf{Instructions.} Read each question carefully to make sure you've answered it fully. You must show your work and explain your reasoning to receive full credit. No calculators or devices of any kind are allowed.

\begin{enumerate}[label=\arabic*), leftmargin=*, itemsep=1em]
    \item Consider the one-dimensional flow defined by
    \[
        u' = r + u - u^3.
    \]
    \begin{enumerate}[label=(\alph*)]
        \item Sketch all the qualitatively different vector fields that occur as the parameter $r$ is varied. Find the value(s) of $r$ at which bifurcations occur, and classify them as saddle-node, transcritical, supercritical pitchfork, or subcritical pitchfork.
        \item Sketch a bifurcation diagram of fixed points $u^*$ as a function of the parameter $r$. Be sure to indicate stability on your diagram.
    \end{enumerate}

    \begin{solution}

    \end{solution}

    \item Consider the one-dimensional flow defined by
    \[
        u' = r(u-1) - u^2(u-1).
    \]
    \begin{enumerate}[label=(\alph*)]
        \item Sketch all the qualitatively different vector fields that occur as the parameter $r$ is varied. Find the value(s) of $r$ at which bifurcations occur, and classify them as saddle-node, transcritical, supercritical pitchfork, or subcritical pitchfork.
        \item Sketch a bifurcation diagram of fixed points $u^*$ as a function of the parameter $r$. Be sure to indicate stability on your diagram.
    \end{enumerate}

    \begin{solution}

    \end{solution}

    \item The following equation describe motion on a line.
    \[
        x' = bx(b-x^2).
    \]
    \begin{enumerate}[label=(\alph*)]
        \item Sketch all the qualitatively different vector fields that occur as $b$ is varied. Find any values of $b$ at which bifurcations occur, and classify these as saddle-node, transcritical, supercritical pitchfork, or subcritical pitchfork.
        \item Sketch the bifurcation diagram of fixed points $x^*$ vs.\ $b$, labeling the stability of each fixed point.
        \item Describe the basin of attraction(s) for the stable fixed point(s) that exist when $b=1$.
    \end{enumerate}

    \begin{solution}

    \end{solution}

    % Source inconsistency: the equation uses r, but all subparts use b.
    \item The following equation describe motion on a line.
    \[
        \frac{dx}{dt} = (r-1+x)(e^x-r).
    \]
    \begin{enumerate}[label=(\alph*)]
        \item Sketch all the qualitatively different vector fields that occur as $b$ is varied. Find any values of $b$ at which bifurcations occur, and classify these as saddle-node, transcritical, supercritical pitchfork, or subcritical pitchfork.
        \item Sketch the bifurcation diagram of fixed points $x^*$ vs.\ $b$, labeling the stability of each fixed point.
        \item Describe the basin of attraction(s) for the stable fixed point(s) that exist when $b=1$.
    \end{enumerate}

    \begin{solution}

    \end{solution}

    \item Consider the system $\dot{x}=ax-x^2+x^3$.
    \begin{enumerate}[label=(\alph*)]
        \item Sketch all the qualitatively different vector fields that occur as $a$ is varied. Find the values of $a$ at which bifurcations occur, and classify these as saddle-node, transcritical, supercritical pitchfork, or subcritical pitchfork.
        \item Sketch the bifurcation diagram of fixed points $x^*$ vs.\ $a$.
        \item Can hysteresis occur in this system? Explain why or why not.
    \end{enumerate}

    \begin{solution}

    \end{solution}

    \item Consider the system $\dot{x}=x-a(x^2+x)$.
    \begin{enumerate}[label=(\alph*)]
        \item Sketch all the qualitatively different vector fields that occur as $a$ is varied. Find the value(s) of $a$ at which bifurcations occur, and classify them as saddle-node, transcritical, supercritical pitchfork, or subcritical pitchfork.
        \item Sketch the bifurcation diagram of fixed points $x^*$ vs.\ $a$.
    \end{enumerate}

    \begin{solution}

    \end{solution}

    \item For the following system of differential equations,
    \[
        \mathbf{x}' = \begin{bmatrix}2 & \alpha^2 \\ -1 & -4\end{bmatrix}\mathbf{x}
    \]
    where $\alpha>0$ is an unspecified constant,
    \begin{enumerate}[label=(\alph*)]
        \item determine the eigenvalues of the coefficient matrix in terms of $\alpha$,
        \item find the critical values of $\alpha$ at which the qualitative nature of the phase portrait experiences a dramatic change, and
        \item sketch the phase portrait for the different choices of $\alpha$: one just below the each critical value and one just above each critical value.
    \end{enumerate}

    \begin{solution}

    \end{solution}

    % Source likely omits a derivative: it prints x = Ax, retained here.
    \item Consider the system 2D linear system $x=Ax$, where
    $A=\begin{bmatrix}a & -1 \\ -1 & b\end{bmatrix}$ for parameters $a>0$ and $b>0$.
    \begin{enumerate}[label=(\alph*)]
        \item Sketch a phase portrait and classify the fixed point at the origin for all possible qualitatively distinct cases as a function of parameters.
        Be sure to include eigenvectors associated with stable/unstable directions in your sketch as appropriate.
        \item Summarize these results by indicating the regions that you've identified in the first quadrant of the $(a,b)$-plane.
    \end{enumerate}

    \begin{solution}

    \end{solution}

    \item Consider the following planar dynamical system:
    \[
        \begin{aligned}
            \dot{x} &= y-y^3 \\
            \dot{y} &= -x-\frac{1}{2}y^2
        \end{aligned}
    \]
    \begin{enumerate}[label=(\alph*)]
        \item Find all fixed points and classify them using eigenvalues of the linearized system. For fixed points with real eigenvalues, find the eigenvectors as well.
        \item Sketch the nullclines and a plausible phase portrait for this system.
    \end{enumerate}

    \begin{solution}

    \end{solution}

    \item Consider the following planar dynamical system:
    \[
        \begin{aligned}
            \dot{x} &= x-xy \\
            \dot{y} &= -y+xy-y^2
        \end{aligned}
    \]
    \begin{enumerate}[label=(\alph*)]
        \item Find all fixed points and classify them using eigenvalues of the linearized system. For fixed points with real eigenvalues, find the eigenvectors as well.
        \item Sketch the nullclines and a plausible phase portrait for this system.
    \end{enumerate}

    \begin{solution}

    \end{solution}

    \item Consider the system
    \[
        \dot{x}=x^2-y, \qquad \dot{y}=x-y.
    \]
    \begin{enumerate}[label=(\alph*)]
        \item Find and classify all fixed points of this system, including eigenvalues and eigenvectors.
        \item Sketch the nullclines and a plausible phase portrait for this system.
    \end{enumerate}

    \begin{solution}

    \end{solution}

    \item Consider the system
    \[
        \dot{x}=x-x^3, \qquad \dot{y}=-y.
    \]
    \begin{enumerate}[label=(\alph*)]
        \item Find and classify all fixed points of this system, including eigenvalues and eigenvectors.
        \item Sketch the nullclines and a plausible phase portrait for this system.
    \end{enumerate}

    \begin{solution}

    \end{solution}

    \item An ODE system models the evolution of social distancing norms during a pandemic.
    The model assumes individuals reduce social contacts when infection rates are high and increase social contacts when infection rates are low. The model has two variables: the level of norm adherence in the community ($x$) and the proportion of infected persons in the community ($y$). The non-dimensionalized equations are shown below:
    \[
        \begin{aligned}
            x' &= ay+b-dx \\
            y' &= r(1-x)y(1-y)-y
        \end{aligned}
    \]
    where $x\in(0,\infty)$ and $y\in[0,1)$. Assume $a\geq 1$, $b>0$, $d>0$ and $r>0$.
    \begin{enumerate}[label=(\alph*)]
        \item Show the model has a unique disease-free equilibrium (DFE) $(x^*,0)$ with $x^*>0$.
        % Source wording retained: part (b)'s positivity claim should be checked
        % against the stated physical domain y in [0,1).
        \item Show there is always at least one endemic equilibrium (EE) $(x^*,y^*)\in\mathbb{R}^2$ with $x^*,y^*>0$
        \par\underline{Hint}: solve for the non-zero equilibrium in $x$. This should yield a quadratic equation of the form $\alpha x^2+\beta x+\gamma=0$. Analyze the coefficients then state conditions for positive roots.
        \item Suppose $b<d$. Show that the DFE is locally stable if $r<r_c$. What is $r_c$?
    \end{enumerate}

    \begin{solution}

    \end{solution}

    \item Consider the 2nd-order ODE that describes the angle $\phi$ of a bead on a rotating hoop:
    \[
        mr\ddot{\phi}+b\dot{\phi}+mg\sin\phi-mr\omega^2\sin\phi\cos\phi=0.
    \]
    Here, $m,b,g,r$ are positive parameters while $\omega\geq 0$.
    \begin{enumerate}[label=(\alph*)]
        \item By nondimensionalizing, show that the ODE can be written in the form
        \[
            \varepsilon\phi''+\phi'+\sin\phi-\gamma\sin\phi\cos\phi=0.
        \]
        for some dimensionless parameters $\varepsilon$ and $\gamma$. Express $\varepsilon$ and $\gamma$ in terms of the parameters of the original system.
        \item Under what assumptions about the original parameters can the system be well-approximated by the following 1D flow?
        \[
            \phi'=-\sin\phi+\gamma\sin\phi\cos\phi
        \]
        \item Assuming that $\gamma=0$, find the fixed points of the above 1D flow and sketch a phase portrait.
        \item Argue that a pitchfork bifurcation occurs at $\gamma=1$. You may use graphical methods or any other valid approach to support your argument.
    \end{enumerate}

    \begin{solution}

    \end{solution}

    % Problems 15--18: M535MT24.pdf (MATH 535 Midterm Exam, Fall 2024).
    \item (25 points) Consider the 1D flow defined by
    \[
        u' = r(u-1)-u^2(u-1).
    \]
    \begin{enumerate}[label=(\alph*)]
        \item Sketch all the qualitatively different vector fields that occur as the parameter $r$ is varied.
        \item Find the value(s) of $r$ at which bifurcations occur, and classify them as saddle-node, transcritical, supercritical pitchfork, or subcritical pitchfork.
        \item Sketch a bifurcation diagram of fixed points $u^*$ as a function of the parameter $r$. Be sure to indicate stability on your diagram.
    \end{enumerate}

    \begin{solution}

    \end{solution}

    \item (25 points) Consider the system
    \[
        \begin{aligned}
            \dot{x} &= -y \\
            \dot{y} &= 2x(1-x^2)
        \end{aligned}
    \]
    \begin{enumerate}[label=(\alph*)]
        \item Find all fixed points of the system.
        \item Classify each of the fixed points you found in part (a), using eigenvalues of the linearized system. For fixed points with real eigenvalues, find the eigenvectors as well.
        \item Show that quantity $E=ax^2+by^2+x^4$ is conserved along trajectories for a suitable choice of $a$ and $b$.
        \item Argue for the existence of periodic orbits in this model. You may use any valid method to support this claim.
        \item Sketch the nullclines and a plausible phase portrait for this system.
    \end{enumerate}

    \begin{solution}

    \end{solution}

    % Source likely omits a derivative: it prints x = Ax, retained here.
    \item (25 points) Consider the system 2D linear system $x=Ax$, where
    $A=\begin{bmatrix}a & -1 \\ -1 & b\end{bmatrix}$ for parameters $a>0$ and $b>0$.
    \begin{enumerate}[label=(\alph*)]
        \item Sketch a phase portrait and classify the fixed point at the origin for all possible qualitatively distinct cases as a function of parameters.
        \par\underline{Note}: Qualitatively accurate phase portraits are sufficient here, it is \emph{not} necessary to compute eigenvectors of $A$.
        \item Summarize these results by indicating the regions that you've identified in the first quadrant of the $(a,b)$-plane.
    \end{enumerate}

    \begin{solution}

    \end{solution}

    \item (25 points) Consider the 2nd-order ODE that describes the angle $\phi$ of a bead on a rotating hoop:
    \[
        mr\ddot{\phi}+b\dot{\phi}+mg\sin\phi-mr\omega^2\sin\phi\cos\phi=0.
    \]
    Here, $m,b,g,r$ are positive parameters while $\omega\geq 0$.
    \begin{enumerate}[label=(\alph*)]
        \item By nondimensionalizing, show that the ODE can be written in the form
        \[
            \varepsilon\phi''+\phi'+\sin\phi-\gamma\sin\phi\cos\phi=0.
        \]
        for some dimensionless parameters $\varepsilon$ and $\gamma$. Express $\varepsilon$ and $\gamma$ in terms of the parameters of the original system.
        \item Under what assumptions about the original parameters can the system be well-approximated by the following 1D flow?
        \[
            \phi'=-\sin\phi+\gamma\sin\phi\cos\phi
        \]
        \item Assuming that $\gamma=0$, find the fixed points of the above 1D flow and sketch a phase portrait.
        \item Argue that a pitchfork bifurcation occurs at $\gamma=1$. You may use graphical methods or any other valid approach to support your argument.
    \end{enumerate}

    \begin{solution}

    \end{solution}
\end{enumerate}

\end{document}